% Iter 64 — Pillar 4 (Length Bias / Dr.GRPO): Conditional Length Response to Reward Direction (CLRRD)
% Per-step (ΔR_t, ΔL_t) trajectory; condition E[ΔL | sign(ΔR), ZVF tier].
% Headline: Dr.GRPO's length growth loses reward-direction coupling, especially
% in the heterogeneous-group (low-ZVF) regime where gradient signal is richest.

\subsubsection{Conditional length response to reward direction}\label{sec:lb-iter64}

Iter~\ref{sec:lb-iter60} separated the marginal-token elasticity
$\epsilon_i = \Delta R_i / \Delta L_i$ into positive- and negative-elasticity
regimes and showed Dr.GRPO spends $\approx 71\%$ of steps in the
destructive regime versus GRPO's $\approx 59\%$. The elasticity analysis
collapsed a 2D trajectory $(\Delta R, \Delta L)$ onto a single ratio.
Iter~\ref{sec:lb-iter64} instead stratifies the same trajectory by the
\emph{sign} of the marginal reward change and asks a sharper question:
when the policy sees a reward-up step, does it differentially compress
or extend its response? We define the \emph{conditional length response}
\begin{equation}\label{eq:lb64-clrrd}
\mathrm{CLRRD}(\mathrm{sign}, \mathrm{ZVF\ tier}) \;=\; \mathbb{E}\!\left[\Delta L_t \mid \mathrm{sign}(\Delta R_t),\ Z_t \in \mathrm{tier}\right]
\end{equation}
and the \emph{reward-direction coupling}
\begin{equation}\label{eq:lb64-alignment}
\mathcal{A} \;=\; \mathbb{E}[\Delta L \mid \Delta R > 0] \;-\; \mathbb{E}[\Delta L \mid \Delta R < 0].
\end{equation}
$\mathcal{A} > 0$ means the policy lengthens on reward-up steps and
shortens on reward-down steps (verbosity-pursuing).
$\mathcal{A} < 0$ means the policy compresses on reward-up steps and
extends on reward-down steps (verbosity-rejecting). Both GRPO and
Dr.GRPO on GSM8K CoT exhibit $\mathcal{A} < 0$ (the policy is
compressing as it learns); the structural question is whether the
\emph{magnitude} of $|\mathcal{A}|$ differs.

\paragraph{Baseline length drift on GSM8K CoT: Dr.GRPO compresses less.}
Pairing the three seeds, $\mathbb{E}[\Delta L]$
(\texttt{length\_bias\_iter64\_summary.tsv}):
\begin{itemize}
\item \textbf{GRPO}: $-0.450$ tokens/step.
\item \textbf{Dr.GRPO}: $-0.241$ tokens/step.
\item \textbf{Paired diff}: $+0.210$, $95\%$ CI $[+0.017,\, +0.405]$,
      $p < 0.001$.
\end{itemize}
Dr.GRPO compresses \emph{less than half} as aggressively as GRPO per
step. This confirms iter56's cumulative-token-tax and iter60's reduced
negative-elasticity fraction: across all iters, Dr.GRPO's per-step
length response is shallower.

\paragraph{Reward-direction coupling collapses in the heterogeneous-group regime.}
On the same three GSM8K CoT seeds, conditioning on ZVF tier
(\texttt{length\_bias\_iter64\_paired.tsv}, cells \texttt{pos\_low},
\texttt{neg\_low}):
\begin{itemize}
\item \textbf{Reward-up, low-ZVF (heterogeneous groups)}:
      GRPO $\mathbb{E}[\Delta L] = -1.84$, Dr.GRPO $= -1.04$,
      paired diff $+0.80$, CI $[+0.12,\, +1.92]$, $p < 0.001$.
      GRPO compresses by $\approx 1.8$ tokens; Dr.GRPO compresses by
      only $\approx 1.0$ tokens -- Dr.GRPO loses $43\%$ of its
      compression signal exactly in the regime where GRPO's gradient
      has the most contrast to work with.
\item \textbf{Reward-up, all tiers}:
      GRPO $= -1.71$, Dr.GRPO $= -1.52$, diff $+0.20$,
      CI $[-0.03,\, +0.36]$, $p = 0.034$. The averaged effect is
      smaller because high-ZVF steps (homogeneous groups) contribute
      nearly identical responses (Dr.GRPO gains nothing on these
      either, but neither does GRPO).
\end{itemize}

\paragraph{Mechanism interpretation.}
The Dr.GRPO estimator is $\hat A_i = \mathrm{clip}\!\left(\frac{r_i - \mu_g}{\sigma_g},\, -c,\, c\right)$
for a prompt-relative baseline $\mu_g$. By clipping the outlier
advantage magnitude, Dr.GRPO equalises the gradient magnitudes
\emph{regardless of the response length}, so the per-token gradient
signal becomes length-invariant. The natural consequence is that
whenever the per-step policy gradient favours short correct responses
(the heterogeneous-group regime), Dr.GRPO under-applies it; whenever
the gradient favours long correct responses (saturation regime), it
over-applies it relative to GRPO. The net result is a more uniform
(less reward-direction-coupled) length trajectory -- the policy is
structurally pulled toward a longer response because the variance-driven
shrinkage signal is damped. Iter64 sharpens the iter60 elasticity
finding by demonstrating that the asymmetry is concentrated in the
low-ZVF regime: GRPO uses \emph{heterogeneous-group contrast} as a
length-compression lever, and Dr.GRPO's clip masks exactly that lever.

\paragraph{Arithmetic-easy task: null finding, by design.}
On the arithmetic-easy task (\texttt{drgrpo\_vs\_grpo.json}, 5 seeds,
40 steps each) all CLRRD coefficients are within $\pm 0.02$ tokens of
zero with CIs that straddle zero. The task converges in $4$--$5$ steps
to reward $\approx 1.0$ with $\mathrm{ZVF} \approx 1.0$ for the
remainder of training, so the trajectory does not visit the
heterogeneous-group regime long enough to develop a length response.
This is a clean negative control: the CLRRD diagnostic is not
spuriously firing on every pair of trajectories.

\paragraph{Reproducibility.}
Driver: \texttt{scripts/length\_bias\_iter64.py} (paired bootstrap,
$2{,}000$ resamples; $3$ GSM8K seeds, $5$ arithmetic seeds; computes
$\mathbb{E}[\Delta L]$ for $12$ (sign $\times$ ZVF) cells per
algorithm). Outputs: \texttt{length\_bias\_iter64\_per\_run.tsv},
\texttt{length\_bias\_iter64\_paired.tsv},
\texttt{length\_bias\_iter64\_summary.tsv}. Figure:
\texttt{figures/length\_bias\_iter64\_clrrd.pdf}, script
\texttt{scripts/length\_bias\_iter64\_fig.py}.